% Part: sets-functions-relations
% Chapter: relations
% Section: relations-as-sets

\documentclass[../../../include/open-logic-section]{subfiles}

\begin{document}

\olfileid{sfr}{rel}{set}
\olsection{संच म्हणून संबंध}

\begin{explain}
\olref[sfr][set][imp]{sec} मध्ये आपण काही महत्त्वाच्या संचांचा उल्लेख केला:
$\Nat$, $\Int$, $\Rat$, $\Real$. यांपैकी काही संचांतील !!{element}s
यांच्यातील काही लक्षवेधी संबंध तुम्हाला आठवत असतील. उदाहरणार्थ, या प्रत्येक
संचावर एक सर्वमान्य \emph{क्रमसंबंध} असतो. तसेच, प्रत्येक वस्तूचा फक्त
स्वतःशी आणि इतर कशाशीही नसणारा \emph{एकरूप असण्याचा} संबंध असतो. पुढे
आपल्याला आणखी अनेक लक्षवेधी संबंध भेटतील, आणि शक्य संबंध तर त्याहूनही
अधिक आहेत. त्यांचा आढावा घेण्यापूर्वी मात्र, संबंधांकडे एका विशिष्ट प्रकारचे
संच म्हणून पाहता येते, हे आपण दाखवू.

यासाठी \olref[sfr][set][pai]{sec} मधील दोन गोष्टी आठवा. प्रथम,
\emph{क्रमित जोडी} ही संकल्पना: $a$ आणि $b$ दिलेले असतील, तर
$\tuple{a, b}$ बनवता येते. येथे घटकांचा क्रम महत्त्वाचा \emph{असतो}.
म्हणून $a \neq b$ असेल, तर $\tuple{a, b} \neq \tuple{b, a}$.
(याची अक्रमित जोड्यांशी, म्हणजे $2$-घटक संचांशी, तुलना करा: तेथे
$\{a, b\}=\{b, a\}$.) दुसरे म्हणजे, \emph{कार्तीय गुणाकार} ही संकल्पना
आठवा: $A$ आणि $B$ हे संच असतील, तर $A \times B$ बनवता येतो. तो
$x \in A$ आणि $y \in B$ असणाऱ्या सर्व $\tuple{x, y}$ जोड्यांचा संच
आहे. विशेषतः, $A^{2}= A \times A$ हा $A$ मधून घेतलेल्या सर्व क्रमित
जोड्यांचा संच आहे.

आता एका संचावरील विशिष्ट संबंध पाहू: नैसर्गिक संख्यांच्या $\Nat$ या
संचावरील $<$-संबंध. $n<m$ असणाऱ्या संख्यांच्या सर्व $\tuple{n, m}$
जोड्यांचा संच विचारात घ्या, म्हणजे,
\[
  R=\Setabs{\tuple{n, m}}{n, m \in \Nat \text{ आणि } n<m}.
\]
$n$ ही संख्या $m$ पेक्षा लहान असणे आणि $\tuple{n, m}$ ही जोडी
$R$ ची सदस्य असणे यांचा निकटचा संबंध आहे, तो असा:
\[
      n<m\text{ तेव्हा आणि केवळ तेव्हाच }\tuple{n, m} \in R.
\]
खरे तर, कोणतीही माहिती न गमावता $R$ हा संच \emph{म्हणजेच} $\Nat$
वरील $<$-संबंध आहे, असे आपण मानू शकतो.

त्याचप्रमाणे, संख्यांमधील कोणत्याही संबंधासाठी $\Nat^{2}$ चा एक उपसंच
तयार करता येतो. उलट, संख्यांच्या जोड्यांचा कोणताही $S \subseteq \Nat^{2}$
संच दिला असेल, तर त्याला अनुरूप संख्यांमधील एक संबंध असतो: $n$ चा $m$
शी तो संबंध तेव्हा आणि केवळ तेव्हाच असतो, जेव्हा $\tuple{n, m} \in S$.
यावरून पुढील व्याख्या समर्थित होते:
\end{explain}

\begin{defn}[द्विपदी संबंध]
$A$ या संचावरील \emph{द्विपदी संबंध} म्हणजे $A^{2}$ चा एक उपसंच.
$R \subseteq A^{2}$ हा $A$ वरील द्विपदी संबंध असेल आणि $x, y \in A$
असतील, तर $\tuple{x, y} \in R$ यासाठी आपण कधी कधी $Rxy$ (किंवा $xRy$)
असे लिहितो.
\end{defn}

\begin{ex}
  \ollabel{relations}
नैसर्गिक संख्यांच्या जोड्यांचा $\Nat^{2}$ हा संच अशा द्विमितीय सारणीत
मांडता येतो:
\[
  \begin{array}{ccccc}
  \mathbf{\tuple{ 0,0 }} & \tuple{ 0,1 } &
    \tuple{ 0,2 } & \tuple{ 0,3 } & \ldots\\
  \tuple{ 1,0 } & \mathbf{\tuple{ 1,1 }} &
    \tuple{ 1,2 } & \tuple{ 1,3 } & \ldots\\
  \tuple{ 2,0 } & \tuple{ 2,1 } &
    \mathbf{\tuple{ 2,2 }} & \tuple{ 2,3 } & \ldots\\
  \tuple{ 3,0 } & \tuple{ 3,1 } & \tuple{ 3,2 } &
    \mathbf{\tuple{ 3,3 }} & \ldots\\
  \vdots & \vdots & \vdots & \vdots & \mathbf{\ddots}
  \end{array}
\]
येथे कर्णावरील नोंदी ठळक केल्या आहेत, कारण त्या जोड्यांचा बनलेला
$\Nat^2$ चा उपसंच, म्हणजे,
\[
  \{\tuple{0,0 }, \tuple{ 1,1 }, \tuple{ 2,2 }, \dots\},
  \]
हा \emph{$\Nat$ वरील एकरूपता संबंध} आहे. (एकरूपता संबंध वारंवार वापरला
जातो, म्हणून कोणत्याही $A$ संचासाठी $\Id{A}=\Setabs{\tuple{ x,x }}{x \in
A}$ अशी व्याख्या करू.) कर्णाच्या वर असलेल्या सर्व जोड्यांचा उपसंच, म्हणजे,
\[
  L = \{\tuple{ 0,1 },\tuple{ 0,2 },\ldots,\tuple{ 1,2 },
  \tuple{ 1,3 }, \dots, \tuple{ 2,3 },\tuple{ 2,4 },\ldots\},
\]
हा \emph{लहान असण्याचा} संबंध आहे, म्हणजे $Lnm$ तेव्हा आणि केवळ तेव्हाच,
जेव्हा $n<m$. कर्णाच्या खाली असलेल्या जोड्यांचा उपसंच, म्हणजे,
\[
  G=\{ \tuple{ 1,0 },\tuple{ 2,0 },\tuple{
    2,1 }, \tuple{ 3,0 },\tuple{ 3,1 },\tuple{ 3,2 }, \dots\},
\]
हा \emph{मोठे असण्याचा} संबंध आहे, म्हणजे $Gnm$ तेव्हा आणि केवळ तेव्हाच,
जेव्हा $n>m$. $L$ आणि $I$ यांचा संयोग $K=L\cup I$ असा लिहू; तो
\emph{लहान किंवा समान असण्याचा} संबंध आहे: $Knm$ तेव्हा आणि केवळ तेव्हाच,
जेव्हा $n \le m$. त्याचप्रमाणे, $H=G \cup I$ हा \emph{मोठे किंवा समान
असण्याचा संबंध} आहे. $L$, $G$, $K$ आणि $H$ हे \emph{क्रम} नावाच्या
विशिष्ट प्रकारचे संबंध आहेत. $L$ आणि $G$ यांचा गुणधर्म असा: कोणत्याही संख्येचा स्वतःशी $L$ किंवा $G$
संबंध नसतो (म्हणजे, सर्व $n$ साठी $Lnn$ आणि $Gnn$ दोन्ही असत्य असतात).
हा गुणधर्म असणाऱ्या संबंधांना \emph{अपरावर्ती} म्हणतात; आणि ते क्रमही
असतील, तर त्यांना \emph{काटेकोर क्रम} म्हणतात.
\end{ex}

\begin{explain}
क्रम आणि एकरूपता हे महत्त्वाचे व स्वाभाविक संबंध असले, तरी आपल्या
व्याख्येनुसार $A^{2}$ चा \emph{कोणताही} उपसंच हा $A$ वरील संबंध असतो,
तो कितीही कृत्रिम किंवा ओढूनताणून बनवलेला वाटला तरीही. विशेषतः,
$\emptyset$ हा कोणत्याही संचावरील संबंध असतो (हा \emph{रिक्त संबंध};
कोणत्याही घटकजोडीमध्ये तो नसतो), आणि $A^{2}$ हाही $A$ वरील संबंध असतो
(प्रत्येक जोडीमध्ये असणारा); त्याला \emph{सार्वत्रिक संबंध} म्हणतात.
तसेच, $E=\Setabs{\tuple{n, m}}{n>5 \text{ किंवा } m \times n \ge 34}$
यासारखी गोष्टही संबंध ठरते.
\end{explain}

\begin{prob}
  $\Pow{\{a, b, c\}}$ या संचावरील $\subseteq$ संबंधाचे !!{element}s लिहा.
\end{prob}

\end{document}
